% !TEX TS-program = lualatex

% ------------------------------------------- %
% -- AUTOMATICALLY GENERATED - DO NOT EDIT -- %
% ------------------------------------------- %

\documentclass{tutodoc}

\usepackage{../preamble.cfg}


\begin{document}

% -- AUTOMATICALLY GENERATED UGLY CODE - START -- %

\section{Supported implementations}
\label{products-all}

All implementations are located in the \tdoccodein{text}+products+ folder. Each implementation provides the following features:

\begin{itemize}
\item
Palette formats in both modular (one file per palette) and monolithic (all palettes in a single file) versions.

\item
Palette definitions in both high-fidelity (original size) and small (currently 40 colors) size.

\item
Select specific colors from an existing palette using their indices.

\item
Shift the palette left (negative value) or right (positive value) by any number of steps.

\item
Reverse the order of the colors.
\end{itemize}

\begin{tdocimp}
To explain how new palettes can be built, we will refer to the colors in the standard palette as \tdoccodein{text}+coul_1+, \tdoccodein{text}+coul_2+, etc., and suppose that the extracted indices are \tdoccodein{text}+{1, 3, 6, 9}+, the shift used is \tdoccodein{text}++1+, and the \tdoccodein{text}+reverse+ option is enabled. The new palette will then be built sequentially as follows: first \tdoccodein{text}+{coul_1, coul_3, coul_6, coul_9}+ (extraction), second \tdoccodein{text}+{coul_9, coul_1, coul_3, coul_6}+ (shift to the right), and finally \tdoccodein{text}+{coul_6, coul_3, coul_1, coul_9}+ (reverse).
\end{tdocimp}

\begin{tdocnote}
Extra features are limited to discrete palette operations. For example, color interpolation is not provided, as this is usually handled out of the box by visualization and formatting tools.
\end{tdocnote}

\subsection{JSON, the versatile default format}
\label{products-json}

The \tdoccodein{text}+JSON+ product allows unsupported programming languages to integrate {\thisproj} palettes easily. Below are the first lines of the \tdoccodein{text}+palettes-hf.json+ file.

\begin{tdoccode}{json}
{
  "Accent": [
    [0.49803922, 0.78823529, 0.49803922],
    [0.74509804, 0.68235294, 0.83137255],
    [0.99215686, 0.75294118, 0.5254902],
    [1, 1, 0.6],
    [0.21960784, 0.42352941, 0.69019608],
    [0.94117647, 0.00784314, 0.49803922],
    [0.74901961, 0.35686275, 0.09019608],
    [0.4, 0.4, 0.4]
  ],
  ...
}
\end{tdoccode}

% -- AUTOMATICALLY GENERATED UGLY CODE - END -- %

\end{document}
